\documentclass{article}
\usepackage[utf8]{inputenc}
\usepackage[T2A]{fontenc}
\usepackage[russian]{babel}
\usepackage{amsmath, amssymb, amsthm}
\usepackage{braket}
\usepackage{geometry}
\geometry{top=2cm, bottom=2cm, left=2.5cm, right=2.5cm}

\newtheorem{theorem}{Теорема}
\newtheorem{definition}{Определение}

\title{Математическая модель GRA-Ramanujan-Generator}
\author{Your Name}
\date{\today}

\begin{document}

\maketitle

\section{Состояния агента}

Каждый агент (машина Рамануджана) описывается набором формул для конечного множества констант $\mathcal{C} = \{c_1,\dots,c_m\}$. Обозначим $\ket{F_{a,c}}$ состояние, кодирующее формулу агента $a$ для константы $c$. Пространство состояний уровня 0 – тензорное произведение по константам:
\[
\ket{\Psi_a^{(0)}} = \bigotimes_{c\in\mathcal{C}} \ket{F_{a,c}} \in \mathcal{H}^{(0)}.
\]

На уровне 1 состояние агента – это просто его уровень 0: $\ket{\Psi_a^{(1)}} = \ket{\Psi_a^{(0)}} \in \mathcal{H}^{(1)} = \mathcal{H}^{(0)}$.

На уровне 2 (популяция) состояние – тензорное произведение состояний всех $N$ агентов:
\[
\ket{\Psi^{(2)}} = \bigotimes_{a=1}^N \ket{\Psi_a^{(1)}} \in \mathcal{H}^{(2)} = \bigotimes_{a=1}^N \mathcal{H}^{(1)}.
\]

\section{Цели и проекторы}

Для каждого уровня определена цель $G_l$ – подпространство «хороших» состояний, и проектор $\mathcal{P}_{G_l}$ на это подпространство.

\subsection{Уровень 0}
Для фиксированной константы $c$ цель $G_{0,c}$ – множество формул с ошибкой $<\varepsilon$ и сложностью $<L$. Проектор $\mathcal{P}_{G_{0,c}}$ можно реализовать как оператор, который упрощает выражение или корректирует коэффициенты, чтобы улучшить точность.

\subsection{Уровень 1}
Цель $G_1$ – согласованность формул для разных констант. Например, если известно тождество $\pi^2/6 = \zeta(2)$, то проектор $\mathcal{P}_{G_1}$ должен сдвигать пару формул $(F_{a,\pi}, F_{a,\zeta(2)})$ к ближайшей паре, удовлетворяющей этому тождеству.

\subsection{Уровень 2}
Цель $G_2$ – оптимальное разнообразие и кооперация в популяции. Проектор $\mathcal{P}_{G_2}$ может обмениваться лучшими формулами между агентами или регулировать веса селекции.

\section{Пена и функционал}

Пена на уровне $l$ – мера отклонения от цели:
\[
\Phi^{(l)}(\Psi^{(l)}, G_l) = \sum_{\mathbf{a}\neq\mathbf{b}} |\bra{\Psi_{\mathbf{a}}^{(l)}} \mathcal{P}_{G_l} \ket{\Psi_{\mathbf{b}}^{(l)}}|^2.
\]

Для практических расчётов мы используем упрощённые выражения:
\begin{align*}
\Phi^{(0)} &= \frac{1}{N}\sum_a \left( \alpha\, \text{error}_a + \beta\, \text{complexity}_a \right),\\
\Phi^{(1)} &= \frac{1}{N}\sum_a \sum_{(c,d)\in I} \bigl( \text{value}(F_{a,c}) - \text{expected}(F_{a,c},F_{a,d}) \bigr)^2,\\
\Phi^{(2)} &= \text{Var}(\{\text{beauty}_a\}).
\end{align*}

Общий функционал:
\[
J = \Lambda_0 \Phi^{(0)} + \Lambda_1 \Phi^{(1)} + \Lambda_2 \Phi^{(2)},
\]
где $\Lambda_l = \lambda_0 \alpha^l$ – веса уровней.

\section{Витальная динамика}

Витальность агента $V_a$ изменяется по закону:
\[
\frac{dV_a}{dt} = \gamma \,\text{beauty}_a - \delta \,\Phi^{(2)}_a - \mu V_a,
\]
где $\Phi^{(2)}_a$ – вклад агента в пену популяции. Агент умирает, если $V_a < V_{\text{death}}$.

\section{Теорема о сходимости}

\begin{theorem}
Если проекторы коммутируют и пространство состояний достаточно велико, то существует состояние $\Psi^*$ с $\Phi^{(l)}(\Psi^*, G_l)=0$ для всех $l$. В этом состоянии все агенты имеют максимальную витальность и их формулы образуют согласованный, разнообразный и точный набор.
\end{theorem}

Доказательство следует индукцией по $l$, аналогично теореме о мультиверсном обнулении.

\end{document}